\section{De-Confounding Ladder Closure}
\label{sec:deconf}

We tested three classes of de-confounding intervention on URTIC + frozen WavLM. All three close negatively at every architectural scope tractable in the project's compute budget. The closure is the negative-result spine of the paper, and each rung's failure mode produced one of the M-disciplines (Section~\ref{ssec:m_disciplines}).

\subsection{A5.5 — Data-level: cross-speaker augmentation}
\label{ssec:a5_5}

\textbf{Audio splicing v1 (killed by M8 self-splice control).} The original A5.5 design spliced anchor chunks with same-cold-different-pseudo-speaker partners using 200 ms crossfades and class-balanced augmentation. A binary splice-detector probe reached UAR 0.998 across every WavLM layer — perfect detection, well above the 0.55 augmentation-class gate. The diagnostic ambiguity ``gate too strict'' vs ``splicer broken'' was resolved by a self-splice control: splicing within the same pseudo-speaker cluster (no cross-speaker mixing) gave UAR 0.990, only $\Delta\,-0.008$ below cross-splice. Cross-speaker mixing accounts for $<$1\,\% of detectability; the splicing operation itself accounts for $\sim$99\,\%. This rules out partner-selection refinements; the splice operation produces a signature regardless of who you splice with. We pivoted to embedding mixup.

\textit{M8 finding (transferable beyond URTIC):} \emph{Audio-level cross-speaker splicing on foundation-model representations is detectable at $>$99\,\% UAR even within a single speaker. The splicing operation itself is the artefact; cross-speaker mixing accounts for $<$1\,\%. Anyone doing audio-level cross-speaker augmentation on FM features should run a self-splice control before training.}

\textbf{Embedding mixup (M9 narrow window empty).} Plan B mixes WavLM pooled stats per chunk: $\tilde{X}_{\text{anchor}} = \alpha \cdot X_{\text{anchor}} + (1-\alpha) \cdot X_{\text{partner}}$, $\alpha \sim U(\alpha_{\min}, \alpha_{\max})$, partner = same-cold + different-pseudo-speaker. K=3 cached mixed variants per chunk, sampled per-epoch by a custom sampler. Two $\alpha$-axis endpoints tested: \emph{conservative} $\alpha \sim U(0.70, 0.85)$ and \emph{aggressive} $\alpha \sim U(0.50, 0.70)$.

\begin{table}[t]
\centering\small
\caption{A5.5 $\alpha$-axis sweep: both endpoints fail differently. M9 narrow-window-empty.}
\label{tab:a55_alpha}
\begin{tabular}{lccr}
\toprule
$\alpha$ range & UAR (3 seeds) & MLP probe & verdict \\
\midrule
A2.5 ref       & $0.6564 \pm 0.0038$ & $0.0501$ & --- \\
$U(0.70, 0.85)$ conservative & $0.6624 \pm 0.0031$ & $0.0506$ & PARTIAL \\
$U(0.50, 0.70)$ aggressive   & $0.6397 \pm 0.0186$ & $0.0485$ & branch (d) FAIL \\
\bottomrule
\end{tabular}
\end{table}

Conservative $\alpha$: small UAR lift ($+0.006$, $\sim$1.6$\sigma$) without measurable probe drop. The augmentation activated weak regularization (recall pattern shifted to more cold-balanced), not de-confounding. cos(init, final) $= 0.9999$ across seeds — the optimizer didn't move the layer weights from A2.5's converged state; the lift is MLP/classifier-level, not layer-weight refinement.

Aggressive $\alpha$: UAR drops by 0.017 ($\sigma$ explodes 6$\times$) AND probes still don't drop. Branch (d) verdict: both endpoints fail differently; intermediate $\alpha \sim U(0.60, 0.80)$ can't escape — gentler falls toward conservative's null, stronger toward aggressive's UAR damage.

\textit{M9 finding:} \emph{Per-chunk pooled-stat mixing on URTIC + frozen WavLM has an empty operating window across the $\alpha$-axis. Run both endpoints before declaring a mixup rung locked. Generalises beyond URTIC: for any cross-speaker augmentation on FM representations operating purely at the data level (without an explicit de-confounding loss), the operating window may be narrow or empty.}

\subsection{A6 — Representation-level head-only: contrastive with speaker-masked positives}
\label{ssec:a6}

A6 adds a fresh projection MLP ($4096 \to 512 \to 128$ with GELU + LayerNorm, L2-normalised on the unit hypersphere) on top of frozen A2.5 standardiser + softmax(layer\_weights). The projection is trained with a supervised contrastive (SupCon) loss with \emph{speaker-masked positives}: same Cold AND different pseudo-speaker. Same-class same-speaker pairs are excluded from positives \emph{and} from the SupCon denominator (the speaker masking lever). Batches are 8 pseudo-speakers $\times$ 8 chunks, 10 epochs, AdamW lr $5\cdot 10^{-5}$, $\tau=0.07$, 3 seeds.

\textbf{Phase 1 PoC.} The PoC reported a striking LR probe drop: $0.0725 \to 0.0410$ ($-43\,\%$ relative, $\sim$13$\sigma$ on the LR seed-$\sigma$). MLP probe barely moved ($0.0501 \to 0.0477$). Class margin emerged ($+0.059$). Cold UAR took a hit ($0.5988$, borderline below the 0.60 floor).

\textbf{M10 — bottleneck-confound diagnostic.} The reflection caught a confound: A2.5's LR probe baseline was on the 4096-d fused representation, while A6's LR probe was on the 128-d L2-normalised z. The 32$\times$ dim reduction plus the unit-hypersphere metric change can drop LR-recoverable speaker info \emph{independent} of the contrastive objective. Three controls on the same 128-d L2-normalised z disambiguate:

\begin{table}[t]
\centering\small
\caption{A6 PoC controls: bottleneck explains the apparent LR drop; cold-CE Pareto-dominates A6 on every metric.}
\label{tab:a6_controls}
\begin{tabular}{lcccc}
\toprule
mode (substrate)        & MLP   & LR    & cold UAR & margin \\
\midrule
A2.5 (4096-d ref)       & $0.0501$ & $0.0725$ & $0.6564$ & --- \\
random proj (128-d)     & $0.0500$ & $0.0427$ & $0.6007$ & $+0.002$ \\
cold-CE only (128-d)    & $0.0408$ & $\mathbf{0.0371}$ & $\mathbf{0.6128}$ & $\mathbf{+0.350}$ \\
vanilla SupCon (128-d)  & $0.0463$ & $0.0400$ & $0.6012$ & $+0.049$ \\
A6 spk-masked SupCon    & $0.0477$ & $0.0410$ & $0.5988$ & $+0.059$ \\
\bottomrule
\end{tabular}
\end{table}

Three independent refutations of the head-only A6 mechanism: (i) random untrained projection drops LR to 0.0427 (within 0.0017 of A6); (ii) cold-CE-only Pareto-dominates A6 on every metric; (iii) vanilla SupCon (no speaker-masking) slightly beats A6 — the speaker-masking lever is mildly counterproductive (reduces positive count without buying de-confounding back).

\textit{M10 finding:} \emph{Probe-substrate dimensionality is a confound for de-confounding measurements. Comparing a probe trained on a high-dim baseline against the same probe on a low-dim post-projection representation conflates information loss with de-confounding. Random-projection controls at the same target dim are mandatory.}

\textbf{M11 — subtractive auxiliary-objective interaction (combined-loss $\lambda$-sweep).} The surviving hypothesis after the controls was ``contrastive as a regulariser on top of CE-anchored cold bottleneck.'' The $\lambda$-sweep $L = L_{\text{cold-CE}} + \lambda \cdot L_{\text{SupCon-spkmasked}}$ over $\lambda \in \{0, 0.05, 0.1, 0.25, 0.5\}$ at 5 seeds gives $\lambda{=}0$ \emph{monotonically} the best on every metric: MLP $0.0408 \to 0.0462$, LR $0.0371 \to 0.0419$, cold UAR $0.6128 \to 0.6005$, class margin $+0.350 \to +0.107$. Adding any positive $\lambda$ strictly worsens every measurement.

\textit{Mechanism (4-step causal chain).} (1) CE on the 128-d projection bottleneck produces sharply class-separated geometry where speaker info is mostly compressed away by the bottleneck itself. (2) Adding SupCon class-pressure flattens within-class diversity — CE already produced separation, contrastive adds no new structure; (3) the flattened diversity includes cold-relevant variation that helped CE generalise — squeezing it out drops cold UAR; (4) SupCon's ``different cold class = separation'' regardless of speaker pulls same-speaker different-cold pairs apart, \emph{re-introducing speaker-correlated variance} and raising the MLP probe.

\textit{M11 finding:} \emph{In low-data settings, an auxiliary contrastive loss is purely subtractive on a CE-anchored bottleneck — when CE has already maximised class margin, contrastive competes with rather than complements that geometry. The naive intuition ``CE for cold + SupCon for de-confounding = both gains stacked'' is empirically wrong here.}

A6 head-only is closed across three recipes: pure speaker-masked SupCon, vanilla SupCon, combined CE+SupCon. Pure cold-CE is the Pareto-best head-only configuration.

\subsection{A7 — Gradient-level: DANN/MDD speaker adversary}
\label{ssec:a7}

A7 inserts a gradient-reversal layer between the projection output $z$ and a pseudo-speaker discriminator. The discriminator learns to predict speaker from $z$ via cross-entropy; the projection sees a $-\lambda$ multiple of the discriminator's gradient and is pushed to make speaker harder to predict (an arms race). Architecture: A2.5 scaler frozen; layer-weights open at lr $10^{-5}$; projection $4096 \to 512 \to 128$ L2-norm fresh; cold head + GRL($\lambda_{\text{adv}}$) $\to$ SpeakerDiscriminator $128 \to 256 \to 210$. $\lambda_{\text{adv}}$ Ganin sigmoid ramp.

\textbf{Phase 1 PoC at 128-d projection.} Sweep $\lambda_{\max} \in \{0, 0.01, 0.03, 0.1, 0.3, 1.0\}$, 3 seeds. Result: \texttt{B\_null} verdict per the discriminator-health diagnostic — discriminator at $\lambda{=}0$ trained to only $\sim$5$\times$ chance (0.0235 vs threshold 0.025 = $5/210$); below this threshold, the verdict is uninterpretable.

\textbf{M12 + M13 — disc-ceiling-before-$\lambda$-sweep ordering.} The disc-ceiling diagnostic loads the saved $\lambda{=}0$ projection checkpoint, freezes it, and trains progressively stronger discriminators on the frozen $z$. The strongest (MLP $128 \to 1024 \to 512 \to 210$, 200 epochs, lr $10^{-3}$) reaches \emph{train top-1 = 0.9521 with devel top-1 = 0.0422 — memorisation gap $+0.910$}. The discriminator memorises individual training-chunk speaker IDs by fingerprint without learning generalising speaker structure that transfers to unseen chunks. The adversary in A7 was therefore pushing against memorised in-batch fingerprints, not generalising speaker direction; no $\lambda > 0$ produced a real probe drop.

\textit{M12 + M13 findings:} \emph{(M12) The discriminator's train/devel memorisation gap is the in-flight validity check during adversarial training: high gap = pushing against memorised fingerprints, no useful adversarial signal. (M13) The disc-ceiling diagnostic must run \emph{before} the $\lambda$-sweep; the decision logic must check \texttt{devel > HIGH} before any memo-gap-as-kill condition. The v1$\to$v2 worked example in this paper (initial fail-fast on memo-gap was over-conservative; corrected logic + regularised disc + lower-$\lambda$ sweep gave \texttt{B\_dann\_dead\_substrate\_resistant}, a stronger closure) demonstrates the two-axis ordering directly.}

\textbf{A7c — un-bottlenecked DANN at 4096-d.} To rule out the projection-bottleneck explanation for A7, we tested DANN directly on the 4096-d fused substrate (no projection MLP at all). Two phases inside one cell, applying M13: (a-i) disc-ceiling pre-flight on frozen 4096-d substrate; (a-ii) adversarial $\lambda$-sweep with regularised discriminator (dropout 0.3, weight\_decay $10^{-3}$).

Pre-flight result: D3 MLP $4096 \to 2048 \to 1024 \to 210$ reaches \emph{train 1.0000 / devel 0.0871, memo gap +0.91}. Devel top-1 (0.084) IS higher at 4096-d than at 128-d (0.042) — the un-bottlenecked substrate retains more recoverable speaker info — but the memorisation gap is identical, so $\sim$91\,\% of apparent speaker recovery at \emph{any} substrate is chunk fingerprinting rather than generalising speaker structure.

After applying the corrected M13 decision (devel $> 0.08$ proceeds regardless of memo gap) and running the $\lambda$-sweep with regularised discriminator: cls UAR, MLP probe, LR probe, AND in-flight disc devel are \emph{all flat across $\lambda \in [0, 0.1]$} — every metric within seed noise of $\lambda{=}0$. \texttt{B\_dann\_dead\_substrate\_resistant}: DANN got a fair fight at every axis the reviewer asked us to check and still didn't shape the substrate.

\textit{M14 finding:} \emph{Speaker-probe-as-de-confounding-measurement has a fundamental noise floor on small-data corpora with coarse pseudo-speaker labels: discriminator memorisation dominates over generalising speaker structure on the substrate REGARDLESS of bottleneck, capping interpretable probe top-1 at $\sim$ceiling $\times$ chance level ($\sim$9\,\% on URTIC at $k{=}210$). Anyone using speaker probes for de-confounding measurement on similar-scale corpora must report disc-ceiling alongside probe, report memorisation gap, and treat probe top-1 as an upper bound on speaker leakage rather than a precise measurement.}

\subsection{Three-level closure summary}

\begin{table}[t]
\centering\small
\caption{De-confounding ladder closure on URTIC + frozen WavLM.}
\label{tab:closure}
\begin{tabular}{llp{3.0cm}}
\toprule
level & rung & verdict \\
\midrule
data         & A5.5 mixup        & M9: empty $\alpha$-window \\
repr (head)  & A6 / A6b          & M10 + M11 across 3 recipes \\
repr (LW)    & A7 PoC            & M10 + M12 mem-dead \\
repr (LW, no proj) & A7c (v2)    & M14: substrate-resistant \\
\bottomrule
\end{tabular}
\end{table}

All within-budget options are exhausted. The methodology framework (M8 -- M14) is the load-bearing contribution; the categorical closure plus the seven instrumentation disciplines are transferable to any small-data paralinguistic foundation-model pipeline.

\subsection{Methodology contributions consolidated}
\label{ssec:m_disciplines}

For convenience, the seven M-disciplines that emerged from the closure are summarised in Table~\ref{tab:m_disciplines}. Each contains a falsifiable mechanism + the failure mode it catches; each generalises beyond URTIC.

\begin{table}[t]
\centering\small
\caption{Seven negative-control disciplines for foundation-model paralinguistic pipelines (M8 -- M14).}
\label{tab:m_disciplines}
\begin{tabular}{lp{4.5cm}}
\toprule
ID & one-line discipline \\
\midrule
M8  & Self-splice control disambiguates ``cross-speaker mixing'' from ``the splicing op itself'' on FM features. \\
M9  & $\alpha$-endpoint controls test both ends of an embedding-mixup design before lock; the operating window may be empty. \\
M10 & Random-projection control at the same target dim is mandatory for any de-confounding rung that introduces a bottleneck; otherwise the apparent probe drop is a bottleneck artefact. \\
M11 & Auxiliary contrastive loss is purely subtractive on a CE-anchored bottleneck; combined-loss $\lambda$-sweep should default to checking $\lambda{=}0$ first. \\
M12 & Discriminator train/devel memorisation gap is the in-flight validity check during adversarial training. \\
M13 & Disc-ceiling diagnostic must run \emph{before} the adversarial $\lambda$-sweep, with $\textit{devel} > \textsc{HIGH}$ checked before any memo-gap-as-kill condition. \\
M14 & Speaker probe top-1 has a memorisation noise floor on small-data corpora with coarse pseudo-speaker labels; treat as upper bound, not precise measurement; report disc-ceiling alongside. \\
\bottomrule
\end{tabular}
\end{table}
